% LOMBEIDA ESCALERAS: Spark es el sistema analizado; RTV es el PFC.
\begin{frame}{10. Spark: recuperar trabajo paralelo}
  \vspace{0.28cm}
  {\large\bfseries Spark reintenta tareas y reconstruye datos perdidos.}\,\cite{sparkConfiguration,sparkRddGuide,sparkFetchFailed}\par
  \vspace{0.60cm}
  \begin{columns}[T,onlytextwidth]
    \begin{column}{0.31\textwidth}
      \begin{block}{Tarea fallida}
        Reintento limitado\,\cite{sparkConfiguration}
      \end{block}
    \end{column}
    \begin{column}{0.31\textwidth}
      \begin{block}{Partición RDD perdida}
        Recalcular por linaje\,\cite{sparkRddGuide}
      \end{block}
    \end{column}
    \begin{column}{0.31\textwidth}
      \begin{block}{Salida shuffle ausente}
        Regenerar datos previos\,\cite{sparkFetchFailed}
      \end{block}
    \end{column}
  \end{columns}
  \vspace{0.48cm}
  {\small\color{textoSuave}La recuperación depende del dato perdido y del tipo de fallo.\,\cite{sparkRddGuide,sparkFetchFailed}}
\end{frame}

\begin{frame}{11. Spark: un fallo real de recuperación}
  {\small\color{textoSuave}SPARK-46182 · registrado en 2023\,\cite{spark46182}}\par
  \vspace{0.55cm}
  \begin{columns}[T,onlytextwidth]
    \begin{column}{0.31\textwidth}
      \begin{block}{1. Tarea termina}
        Crea nuevos bloques shuffle.\,\cite{spark46182}
      \end{block}
    \end{column}
    \begin{column}{0.31\textwidth}
      \begin{block}{2. Executor sale}
        La migración aún no los detecta.\,\cite{spark46182}
      \end{block}
    \end{column}
    \begin{column}{0.31\textwidth}
      \begin{block}{3. Datos perdidos}
        Se recalcula la partición.\,\cite{spark46182}
      \end{block}
    \end{column}
  \end{columns}
  \vspace{0.55cm}
  {\small\bfseries Causa raíz:} carrera entre el fin de la tarea y el retiro del executor.\,\cite{spark46182}
\end{frame}

\begin{frame}{12. Lecciones para RTV y conclusiones}
  \vspace{0.28cm}
  {\small\color{textoSuave}Propuestas de mejora para el RTV; no se afirma que ya estén implementadas.}\par
  \vspace{0.55cm}
  \begin{columns}[T,onlytextwidth]
    \begin{column}{0.47\textwidth}
      \begin{block}{1. Conservar el avance}
        Confirmar cada resultado antes de continuar.\,\cite{postgresTransactions}
      \end{block}
    \end{column}
    \begin{column}{0.47\textwidth}
      \begin{block}{2. Cerrar sin perder datos}
        Verificar el estado completo antes de certificar.\,\cite{spark46182,postgresTransactions}
      \end{block}
    \end{column}
  \end{columns}
  \vspace{0.60cm}
  {\small\bfseries Idea final:} atender vehículos a la vez exige conservar íntegro el resultado de cada uno.\,\cite{spark46182,postgresTransactions}
\end{frame}

